\documentclass[12pt,a4paper,landscape]{article}
\usepackage[utf8]{inputenc}
\usepackage[T1]{fontenc}
\usepackage[french]{babel}
\usepackage[margin=1.5cm]{geometry}
\usepackage{xcolor}
\usepackage{booktabs}
\usepackage{longtable}
\usepackage{array}
\usepackage{colortbl}
\usepackage{fancyhdr}
\usepackage{hyperref}
\usepackage{tabularx}
\usepackage{multirow}

% ---- Couleurs ----
\definecolor{headergreen}{RGB}{92, 184, 92}
\definecolor{headerdark}{RGB}{60, 118, 61}
\definecolor{titlebg}{RGB}{70, 180, 210}
\definecolor{rowtitle}{RGB}{220, 53, 69}
\definecolor{rowalt}{RGB}{245, 255, 245}
\definecolor{passgreen}{RGB}{40, 167, 69}
\definecolor{failred}{RGB}{220, 53, 69}
\definecolor{lightgray}{RGB}{248, 248, 248}

\pagestyle{fancy}
\fancyhf{}
\fancyhead[L]{\textbf{Devoir 304 — Cas de test — Systeme de transaction bancaire}}
\fancyhead[R]{Avril 2026}
\fancyfoot[C]{\thepage}
\renewcommand{\headrulewidth}{0.5pt}

\setlength{\parindent}{0pt}
\setlength{\parskip}{6pt}

% Colonne à largeur fixe avec retour à la ligne
\newcolumntype{L}[1]{>{\raggedright\arraybackslash}p{#1}}

\begin{document}

% =========================================================
% TITRE
% =========================================================
\begin{center}
  {\Large\bfseries\colorbox{titlebg}{\textcolor{white}{\quad Catégories de tests logiciels \quad}}}\\[0.5cm]
  {\large\bfseries\textcolor{rowtitle}{Conception et gestion des cas de test}}\\[0.2cm]
  {\large\bfseries\textcolor{blue}{Format de cas de test (modèle 1)}}\\[0.5cm]
  {\normalsize API de transaction bancaire — 6 fonctions — 29 cas de test}
\end{center}

% =========================================================
% FONCTION 1 — Créer un compte
% =========================================================
\vspace{0.3cm}
{\bfseries\large\textcolor{rowtitle}{Fonction 1 : Créer un compte \quad \texttt{POST /accounts}}}
\vspace{0.2cm}

\begin{longtable}{|L{2cm}|L{5cm}|L{6.5cm}|L{5cm}|L{2.5cm}|}
\hline
\rowcolor{headergreen}
\textbf{\textcolor{white}{Test Case ID}} &
\textbf{\textcolor{white}{Test Case Description}} &
\textbf{\textcolor{white}{Test Input (Test Case Value)}} &
\textbf{\textcolor{white}{Expected Results}} &
\textbf{\textcolor{white}{Test Case Status}} \\
\hline
CT-01 &
Création valide avec tous les champs &
\texttt{full\_name: "KOA ESSAMA PAULIN BRICE"} \newline \texttt{phone\_number: "699001122"} \newline \texttt{email: "paulin@example.com"} \newline \texttt{initial\_balance: 10000} &
HTTP 201 — Compte créé avec un \texttt{id} UUID et un \texttt{account\_number} générés automatiquement. \texttt{balance = 10000} &
PASS \\
\hline
\rowcolor{rowalt}
CT-02 &
Création valide sans email ni solde initial &
\texttt{full\_name: "NGUEMA CAROLE"} \newline \texttt{phone\_number: "677445566"} &
HTTP 201 — Compte créé avec \texttt{balance = 0} et \texttt{email = null} &
PASS \\
\hline
CT-03 &
Champ \texttt{full\_name} absent &
\texttt{phone\_number: "699001122"} \newline \texttt{initial\_balance: 5000} &
HTTP 400 / 422 — Erreur : \textit{"full\_name et phone\_number sont obligatoires."} &
PASS \\
\hline
\rowcolor{rowalt}
CT-04 &
Champ \texttt{phone\_number} absent &
\texttt{full\_name: "KOA ESSAMA PAULIN BRICE"} \newline \texttt{initial\_balance: 5000} &
HTTP 400 / 422 — Erreur : champ obligatoire manquant &
PASS \\
\hline
CT-05 &
Solde initial négatif &
\texttt{full\_name: "TEST"} \newline \texttt{phone\_number: "600000001"} \newline \texttt{initial\_balance: -500} &
HTTP 400 / 422 — Erreur : \textit{"Le solde initial doit être positif ou nul."} &
PASS \\
\hline
\end{longtable}

% =========================================================
% FONCTION 2 — Lister les comptes
% =========================================================
\vspace{0.3cm}
{\bfseries\large\textcolor{rowtitle}{Fonction 2 : Lister les comptes \quad \texttt{GET /accounts}}}
\vspace{0.2cm}

\begin{longtable}{|L{2cm}|L{5cm}|L{6.5cm}|L{5cm}|L{2.5cm}|}
\hline
\rowcolor{headergreen}
\textbf{\textcolor{white}{Test Case ID}} &
\textbf{\textcolor{white}{Test Case Description}} &
\textbf{\textcolor{white}{Test Input (Test Case Value)}} &
\textbf{\textcolor{white}{Expected Results}} &
\textbf{\textcolor{white}{Test Case Status}} \\
\hline
CT-06 &
Liste vide au démarrage &
Aucun compte créé. \newline Requête : \texttt{GET /accounts} &
HTTP 200 — Corps : \texttt{[]} (tableau vide) &
PASS \\
\hline
\rowcolor{rowalt}
CT-07 &
Liste avec comptes existants &
Au moins un compte créé au préalable. \newline Requête : \texttt{GET /accounts} &
HTTP 200 — Tableau de comptes sans la liste des transactions &
PASS \\
\hline
\end{longtable}

% =========================================================
% FONCTION 3 — Consulter un compte
% =========================================================
\vspace{0.3cm}
{\bfseries\large\textcolor{rowtitle}{Fonction 3 : Consulter un compte \quad \texttt{GET /accounts/\{account\_id\}}}}
\vspace{0.2cm}

\begin{longtable}{|L{2cm}|L{5cm}|L{6.5cm}|L{5cm}|L{2.5cm}|}
\hline
\rowcolor{headergreen}
\textbf{\textcolor{white}{Test Case ID}} &
\textbf{\textcolor{white}{Test Case Description}} &
\textbf{\textcolor{white}{Test Input (Test Case Value)}} &
\textbf{\textcolor{white}{Expected Results}} &
\textbf{\textcolor{white}{Test Case Status}} \\
\hline
CT-08 &
Consulter un compte existant &
\texttt{account\_id} = UUID valide récupéré après création &
HTTP 200 — Détails du compte avec l'historique complet des transactions &
PASS \\
\hline
\rowcolor{rowalt}
CT-09 &
Consulter un compte inexistant &
\texttt{account\_id} = \texttt{00000000-0000-0000-0000-000000000000} (ID invalide) &
HTTP 404 — \textit{"Compte introuvable."} &
PASS \\
\hline
\end{longtable}

% =========================================================
% FONCTION 4 — Dépôt
% =========================================================
\vspace{0.3cm}
{\bfseries\large\textcolor{rowtitle}{Fonction 4 : Effectuer un dépôt \quad \texttt{POST /accounts/\{account\_id\}/deposit}}}
\vspace{0.2cm}

\begin{longtable}{|L{2cm}|L{5cm}|L{6.5cm}|L{5cm}|L{2.5cm}|}
\hline
\rowcolor{headergreen}
\textbf{\textcolor{white}{Test Case ID}} &
\textbf{\textcolor{white}{Test Case Description}} &
\textbf{\textcolor{white}{Test Input (Test Case Value)}} &
\textbf{\textcolor{white}{Expected Results}} &
\textbf{\textcolor{white}{Test Case Status}} \\
\hline
CT-10 &
Dépôt valide sans description &
Compte avec \texttt{balance = 10000} \newline \texttt{amount: 5000} &
HTTP 200 — \texttt{balance = 15000}. Transaction de type \texttt{deposit} enregistrée &
PASS \\
\hline
\rowcolor{rowalt}
CT-11 &
Dépôt valide avec description &
Compte existant \newline \texttt{amount: 3000} \newline \texttt{description: "Depot agence"} &
HTTP 200 — Solde augmente, libellé visible dans l'historique &
PASS \\
\hline
CT-12 &
Montant égal à zéro &
Compte existant \newline \texttt{amount: 0} &
HTTP 400 / 422 — \textit{"Le montant doit être strictement positif."} &
PASS \\
\hline
\rowcolor{rowalt}
CT-13 &
Montant négatif &
Compte existant \newline \texttt{amount: -1000} &
HTTP 400 / 422 — \textit{"Le montant doit être strictement positif."} &
PASS \\
\hline
CT-14 &
Dépôt sur un compte inexistant &
\texttt{account\_id} invalide \newline \texttt{amount: 5000} &
HTTP 404 — \textit{"Compte introuvable."} &
PASS \\
\hline
\end{longtable}

% =========================================================
% FONCTION 5 — Retrait
% =========================================================
\vspace{0.3cm}
{\bfseries\large\textcolor{rowtitle}{Fonction 5 : Effectuer un retrait \quad \texttt{POST /accounts/\{account\_id\}/withdraw}}}
\vspace{0.2cm}

\begin{longtable}{|L{2cm}|L{5cm}|L{6.5cm}|L{5cm}|L{2.5cm}|}
\hline
\rowcolor{headergreen}
\textbf{\textcolor{white}{Test Case ID}} &
\textbf{\textcolor{white}{Test Case Description}} &
\textbf{\textcolor{white}{Test Input (Test Case Value)}} &
\textbf{\textcolor{white}{Expected Results}} &
\textbf{\textcolor{white}{Test Case Status}} \\
\hline
CT-15 &
Retrait valide, solde suffisant &
Compte avec \texttt{balance = 10000} \newline \texttt{amount: 2000} \newline \texttt{description: "Retrait DAB"} &
HTTP 200 — \texttt{balance = 8000}. Transaction \texttt{withdraw} enregistrée &
PASS \\
\hline
\rowcolor{rowalt}
CT-16 &
Retrait égal au solde disponible &
Compte avec \texttt{balance = 10000} \newline \texttt{amount: 10000} &
HTTP 200 — \texttt{balance = 0}. Compte vidé, transaction enregistrée &
PASS \\
\hline
CT-17 &
Solde insuffisant &
Compte avec \texttt{balance = 1000} \newline \texttt{amount: 50000} &
HTTP 400 — \textit{"Solde insuffisant pour effectuer ce retrait."} &
PASS \\
\hline
\rowcolor{rowalt}
CT-18 &
Montant égal à zéro &
Compte existant \newline \texttt{amount: 0} &
HTTP 400 / 422 — \textit{"Le montant doit être strictement positif."} &
PASS \\
\hline
CT-19 &
Montant négatif &
Compte existant \newline \texttt{amount: -500} &
HTTP 400 / 422 — \textit{"Le montant doit être strictement positif."} &
PASS \\
\hline
\rowcolor{rowalt}
CT-20 &
Retrait sur un compte inexistant &
\texttt{account\_id} invalide \newline \texttt{amount: 1000} &
HTTP 404 — \textit{"Compte introuvable."} &
PASS \\
\hline
\end{longtable}

% =========================================================
% FONCTION 6 — Virement
% =========================================================
\vspace{0.3cm}
{\bfseries\large\textcolor{rowtitle}{Fonction 6 : Virement entre comptes \quad \texttt{POST /accounts/\{account\_id\}/transfer}}}
\vspace{0.2cm}

\begin{longtable}{|L{2cm}|L{5cm}|L{6.5cm}|L{5cm}|L{2.5cm}|}
\hline
\rowcolor{headergreen}
\textbf{\textcolor{white}{Test Case ID}} &
\textbf{\textcolor{white}{Test Case Description}} &
\textbf{\textcolor{white}{Test Input (Test Case Value)}} &
\textbf{\textcolor{white}{Expected Results}} &
\textbf{\textcolor{white}{Test Case Status}} \\
\hline
CT-24 &
Virement valide entre deux comptes &
Compte A : \texttt{balance = 10000} \newline Compte B : \texttt{balance = 5000} \newline \texttt{to\_account\_id} = ID du compte B \newline \texttt{amount: 3000} &
HTTP 200 — Compte A : \texttt{balance = 7000}, Compte B : \texttt{balance = 8000}. \newline Transactions \texttt{transfer\_out} (A) et \texttt{transfer\_in} (B) enregistrées &
PASS \\
\hline
\rowcolor{rowalt}
CT-25 &
Virement valide avec description &
Compte A et B existants \newline \texttt{amount: 1000} \newline \texttt{description: "Remboursement loyer"} &
HTTP 200 — Description visible dans les deux historiques de transactions &
PASS \\
\hline
CT-26 &
Solde insuffisant pour le virement &
Compte A : \texttt{balance = 1000} \newline \texttt{amount: 50000} &
HTTP 400 — \textit{"Solde insuffisant pour effectuer ce virement."} &
PASS \\
\hline
\rowcolor{rowalt}
CT-27 &
Virement vers le même compte &
\texttt{to\_account\_id} = même \texttt{account\_id} que la source \newline \texttt{amount: 1000} &
HTTP 400 — \textit{"Le compte source et le compte destinataire doivent être différents."} &
PASS \\
\hline
CT-28 &
Compte source inexistant &
\texttt{account\_id} invalide \newline \texttt{to\_account\_id} valide \newline \texttt{amount: 1000} &
HTTP 404 — \textit{"Compte introuvable."} &
PASS \\
\hline
\rowcolor{rowalt}
CT-29 &
Compte destinataire inexistant &
\texttt{account\_id} valide \newline \texttt{to\_account\_id} = \texttt{00000000-0000-0000-0000-000000000000} \newline \texttt{amount: 1000} &
HTTP 404 — \textit{"Compte destinataire introuvable."} &
PASS \\
\hline
\end{longtable}

% =========================================================
% FONCTION + — Historique des transactions
% =========================================================
\vspace{0.3cm}
{\bfseries\large\textcolor{rowtitle}{Fonction + : Historique des transactions \quad \texttt{GET /accounts/\{account\_id\}/transactions}}}
\vspace{0.2cm}

\begin{longtable}{|L{2cm}|L{5cm}|L{6.5cm}|L{5cm}|L{2.5cm}|}
\hline
\rowcolor{headergreen}
\textbf{\textcolor{white}{Test Case ID}} &
\textbf{\textcolor{white}{Test Case Description}} &
\textbf{\textcolor{white}{Test Input (Test Case Value)}} &
\textbf{\textcolor{white}{Expected Results}} &
\textbf{\textcolor{white}{Test Case Status}} \\
\hline
CT-21 &
Historique d'un compte sans opération &
Compte créé avec \texttt{initial\_balance = 0} \newline \texttt{GET /accounts/\{id\}/transactions} &
HTTP 200 — Corps : \texttt{[]} (aucune transaction) &
PASS \\
\hline
\rowcolor{rowalt}
CT-22 &
Historique avec plusieurs opérations &
Compte avec dépôt, retrait et virement effectués \newline \texttt{GET /accounts/\{id\}/transactions} &
HTTP 200 — Liste chronologique de toutes les transactions avec \texttt{transaction\_type}, \texttt{amount}, \texttt{balance\_after} &
PASS \\
\hline
CT-23 &
Historique d'un compte inexistant &
\texttt{account\_id} invalide &
HTTP 404 — \textit{"Compte introuvable."} &
PASS \\
\hline
\end{longtable}

\vfill
\begin{center}
\rule{0.5\linewidth}{0.5pt}\\[0.2cm]
{\small Total : \textbf{29 cas de test} couvrant les 6 fonctions du système}
\end{center}

\end{document}
